\documentclass[11pt]{letter}
\usepackage[a4paper,margin=1in]{geometry}
\usepackage{mathptmx}
\usepackage{microtype}
\newcommand{\dd}{\mathrm{d}}

\signature{Daryl Janzen}
\address{Department of Physics and Engineering Physics\\
University of Saskatchewan\\
116 Science Place\\
Saskatoon, SK S7N 5E2, Canada\\
daryl.janzen@usask.ca}
\date{\today}

\begin{document}
\begin{letter}{The Editors\\
\textit{Classical and Quantum Gravity}}

\opening{Dear Editors,}

I submit for consideration as a Paper the manuscript ``The Schwarzschild horizon and singularity as the two poles of one circle: the analytic structure of the maximally extended geometry.''

The event horizon and the curvature singularity of the Schwarzschild geometry are usually classified as different in kind: the one a coordinate artefact that a better chart removes, the other an edge at which the description must stop. The manuscript examines that classification at the level of the analytic structure. In the Lema\^itre--Tolman cycloid parametrisation of radial infall, $r(z)=M(1+\cos z)$, the horizon and $r=0$ are the two non-degenerate critical points of one real-analytic curve, the poles of the circle $(r-M)^2+(\dd r/\dd z)^2=M^2$. The same operation, continuation of the parameter in an imaginary direction, carries the curve through either: at the horizon to the exterior regions of the Kruskal--Szekeres extension, and at $r=0$ to two regions isometric to Schwarzschild of negative mass. The manuscript proves these are the only real sections of the continued curve.

What separates the two loci is the value of the areal radius at the critical point, and the manuscript shows that every invariant difference between them, including the order-12 pole of the Kretschmann scalar and its approach to the Kasner value $Ks^4\to64/27$, is carried by that one number. None of this contradicts the established results: the curvature singularity is genuine, Sbierski's $C^0$-inextendibility theorem stands, and nothing physical is said to pass through $r=0$. The claim concerns what the classification may be read as saying, not which spacetime is the maximal one.

I believe the paper belongs in \textit{Classical and Quantum Gravity} because it bears on how singularities of exact solutions are classified and on what the analytic extension of a solution does and does not establish, questions this journal has long carried. Every computation in the paper has been checked symbolically, and the scripts are available on request.

The manuscript has not been published and is not under consideration elsewhere.

\closing{Sincerely,}

\end{letter}
\end{document}
